% Fichier complété par tools/complete_missing_corrections.py.
% Source : SLCI/SLCI-03-SchemaBlocs/96_Stabilisateur/96_Stabilisateur.tex
% Statut : rédigé (1 question(s)).
\begin{corrigebox}[Corrigé rédigé]
\CorrigeQuestion{1}
On déplace successivement les sommateurs et les points de prélèvement
        en appliquant les règles de transposition : lorsqu'un sommateur traverse un bloc
        $G$, la branche déplacée est multipliée ou divisée par $G$ afin de conserver le
        même signal. Puis on regroupe les blocs en série par produit, les blocs en
        parallèle par somme et les boucles simples par $G/(1+GH)$. Les deux schémas ont
        alors les mêmes relations algébriques entrée--sortie et perturbation--sortie.
\end{corrigebox}
